\documentclass[11pt]{article}
\usepackage{amsmath,amssymb,booktabs}
\usepackage[margin=1in]{geometry}

\newcommand{\Psym}{\mathsf P_\infty}
\newcommand{\Mz}{M_\zeta}

\title{Step 175: Numerical Evaluation of \(L_{\rho_1,0}(G_*)\)}
\author{Six Birds Foundations III}
\date{2026-05-15}

\begin{document}
\maketitle

\section{Verdict}

\[
\boxed{\texttt{L\_verdict = V\_L\_numerical\_nonzero}.}
\]

Under the Step 173/174 operational identity and the declared \(\lambda=1\)
normalization,
\[
L_{\rho_1,0}(G_*)
=-I_{\rm sinc}-\sum_{n\ge0}\Psi_n(\rho_1)J_n,
\]
with
\[
\rho_1=\frac12+14.134725141734693\,i.
\]
The numerical evaluation with the required \(n=0,1,2\) PSWF block gives
\[
\boxed{
L_{\rho_1,0}^{[0,2]}(G_*)
=0.12146653476134234
-0.08985731971133760\,i
}
\]
with total error radius
\[
\boxed{1.9648066412630187\times 10^{-4}.}
\]
Therefore
\[
\boxed{|L_{\rho_1,0}(G_*)|\ge 0.1508944091096583}
\]
within the numerical error budget used here.  The value is separated from zero.

\section{T1. Moments and Coefficients}

The Step 174 generator is
\[
g_*(t)=b_1(t)-\alpha b_2(t)+\beta_3 b_3(t),
\quad
b_j(t)=\beta\!\left(\frac{t-c_j}{1/5}\right),
\]
with \(c_1=3/2\), \(c_2=5/2\), \(c_3=7/2\), and
\[
\beta(u)=e^{-1/(1-u^2)}\quad (|u|<1).
\]
High-order Gauss--Legendre/scipy quadrature gives:
\[
\begin{array}{c|r}
\toprule
\text{quantity} & \text{value}\\
\midrule
A_1 & 8.8798763233615893\cdot10^{-2}\\
A_2 & 8.8798763233615963\cdot10^{-2}\\
A_3 & 8.8798763233615963\cdot10^{-2}\\
B_1 & 5.9366578012354863\cdot10^{-2}\\
B_2 & 3.5555525747111349\cdot10^{-2}\\
B_3 & 2.5384188380051690\cdot10^{-2}\\
\alpha & 3.3409952306131059\\
\beta_3 & 2.3409952306131068\\
\bottomrule
\end{array}
\]
The numerical endpoint checks are
\[
G_*(0)=8.3266726846886741\cdot10^{-17},
\qquad
G_*(1)=5.5511151231257827\cdot10^{-17},
\]
consistent with the Step 119 legality conditions.

\section{T2--T3. Sampling \(G_*\) and \(\zeta\)}

For \(s=1/2+iu\),
\[
G_*(s)=\int_1^4 g_*(t)t^{-s}\,dt
=\int_1^4 g_*(t)t^{-1/2}e^{-iu\log t}\,dt.
\]
The computation used Gauss--Legendre quadrature on each bump support with
\(280\) nodes per support, hence \(840\) \(t\)-nodes total.  The zeta values
\(\zeta(1/2+iu)\) were computed by \texttt{mpmath.zeta} at \(50\) decimal
digits on the grid
\[
u\in[-200,200],\qquad h=0.05.
\]
Selected sampled values are recorded in
\texttt{G\_zeta\_samples\_step175.csv}.

\section{T4. Sinc Integral}

The sinc integral is
\[
I_{\rm sinc}
=\int_{\mathbb R}
\frac{\sin(\gamma_1-u)}{\pi(\gamma_1-u)}
\zeta(1/2+iu)G_*(1/2+iu)\,du,
\]
where the removable singularity uses the limiting value \(1/\pi\).
Numerically,
\[
\boxed{
I_{\rm sinc}
=-0.11261872344887212
+0.08976645970797142\,i
}
\]
and
\[
|I_{\rm sinc}|=0.14401733978850045.
\]
The error radius assigned to this component is
\[
1.8185144104892716\cdot10^{-4},
\]
composed of:
\[
\begin{array}{c|r}
\toprule
\text{component} & \text{size}\\
\midrule
\text{grid refinement } h=0.10\to0.05 & 6.4362848920259704\cdot10^{-10}\\
\text{tail shell } [160,200] \text{ doubled} & 1.8185079641813977\cdot10^{-4}\\
\text{\(G_*\) quadrature comparison } 180\to280 & 2.2981884777988851\cdot10^{-15}\\
\text{zeta/double conversion allowance} & 10^{-12}\\
\bottomrule
\end{array}
\]

\section{T5. PSWF Sum}

The Step 173 basis was computed by diagonalizing the sinc kernel
\[
K(x,y)=\frac{\sin(x-y)}{\pi(x-y)},\qquad x,y\in[-1,1],
\]
with \(320\) Gauss--Legendre nodes.  This implements the same normalization as
the operational identity:
\[
\Psi_n=\mathcal U_\infty
\left(\frac{(I-P_1)\phi_n^1}{\sqrt{1-\mu_n^1}}\right).
\]
For \(n=0,1,2\):
\[
\begin{array}{c|r|r|r|r}
\toprule
n & \mu_n & \Psi_n(\rho_1) & J_n & \Psi_n(\rho_1)J_n\\
\midrule
0 & 5.7258178063789256\cdot10^{-1}
& 4.1322189808145601\cdot10^{-2}
& -2.1260953610400818\cdot10^{-1}
& -8.7854916059116103\cdot10^{-3}\\
1 & 6.2791274149805884\cdot10^{-2}
& -9.2849186138664966\cdot10^{-4}
& -9.7857619592373851\cdot10^{-2}i
& 9.0860003366189876\cdot10^{-5}i\\
2 & 1.2374793284661086\cdot10^{-3}
& 3.0995624623184163\cdot10^{-3}
& -2.0105968928270435\cdot10^{-2}
& -6.2319706558607477\cdot10^{-5}\\
\bottomrule
\end{array}
\]
Thus
\[
\sum_{n=0}^{2}\Psi_n(\rho_1)J_n
=-0.008847811312470217
+0.00009086000336618477\,i.
\]
The extended diagnostic through \(n=11\) gives
\[
\sum_{n=0}^{11}\Psi_n(\rho_1)J_n
=-0.008847825468795223
+0.00009125082134558793\,i.
\]
The observed absolute sum of \(n=3,\ldots,11\) terms is
\[
4.070701178842979\cdot10^{-7},
\]
and the conservative truncation bound used for \(n\ge3\) is
\[
1.0\cdot10^{-6}.
\]

\section{T6. Final Numerical Value and Error Budget}

Using the required \(n=0,1,2\) block,
\[
\begin{aligned}
L_{\rho_1,0}^{[0,2]}(G_*)
&=-I_{\rm sinc}
-\sum_{n=0}^{2}\Psi_n(\rho_1)J_n\\
&=
0.12146653476134234
-0.08985731971133760\,i.
\end{aligned}
\]
The total error radius is
\[
\begin{array}{c|r}
\toprule
\text{component} & \text{radius}\\
\midrule
I_{\rm sinc}\ \text{quadrature/tail} & 1.8185144104892716\cdot10^{-4}\\
\text{PSWF \(n=0,1,2\) grid/normalization diagnostic} & 1.3629222077374727\cdot10^{-5}\\
\text{PSWF tail bound \(n\ge3\)} & 1.0\cdot10^{-6}\\
\midrule
\text{total} & 1.9648066412630187\cdot10^{-4}\\
\bottomrule
\end{array}
\]
Since
\[
|L_{\rho_1,0}^{[0,2]}(G_*)|
=0.15109088977378460,
\]
we obtain the lower bound
\[
0.15109088977378460-0.00019648066412630187
=0.1508944091096583.
\]

\section{Branch C Status}

This gives numerical nonzero evidence for the selected nontrivial legal
generator.  Under the Step 173/174 normalization, the full-carrier annihilation
route is not supported at \((\rho_1,0,G_*)\).  Branch C therefore requires
strict subclass restriction, a normalization audit, or an external analytic
countercheck before any global conclusion can be drawn.

\section{Nonclaim Boundary}

This step does not prove RH, does not prove \(\Xi_{\mathrm{BC}}=0\), does not
close Branch C globally, and does not remove any retained no-go.  It is a
numerical matrix-element evaluation for one generator and one zero label.

\end{document}
